\chapter{Introdução e Fundamentação}

\section{O Problema da Degradação Conceitual em Agentes Autônomos}
A utilização de modelos de linguagem de grande porte (LLMs) em tarefas de desenvolvimento de software trouxe capacidade de escrita de código em alta velocidade. Contudo, em projetos extensos, a atuação não estruturada de múltiplos agentes acarreta o fenômeno da \emph{crise conceitual} ou \emph{context drift}. Quando agentes especialistas atuam sem um plano unificado, decisões locais violam invariantes de arquitetura globais.

\section{A Filosofia Quantiflux}
O \textbf{Quantiflux} foi concebido para impor disciplina estrutural estrita antes da geração de qualquer linha de código. A premissa fundamental do framework é que a janela de contexto de um agente de nível estratégico é um ativo de alto valor. Para preservá-la, a execução de código é delegada a subagentes táticos temporários e isolados, operando sob o controle determinístico de um Plano Diretor de Engenharia (\emph{Master Engineering Plan}).

\section{Objetivos da Metodologia}
\begin{itemize}
    \item Formalizar a separação de responsabilidades em quatro níveis hierárquicos.
    \item Estabelecer a execução sequencial por Ondas (\emph{Waves}) com portões de qualidade invioláveis.
    \item Garantir que mudanças estruturais drásticas sejam tratadas no nível de tese e plano diretor, e não como improvisos operacionais.
\end{itemize}